Steering vectors---directions in activation space that causally shift model behavior---are an increasingly popular tool for controllable generation and AI safety. They rest on the assumption that concepts are represented as linear directions, a view formalized by the Linear Representation Hypothesis. But does linear \emph{detectability} imply linear \emph{steerability}? We investigate this question for 11 emotional styles in \pythia by comparing three measures: (i)~linear versus MLP probe accuracy, (ii)~contrastive activation addition effectiveness, and (iii)~prompting success. We find that all 11 styles are represented linearly (non-linearity indices 0.965--1.008), yet linear separability does not predict steering effectiveness ($r = -0.220$, $p = 0.52$). Prompting succeeds uniformly for all styles (3.8--4.4/5 adherence), including those that steer poorly, and is uncorrelated with steering success ($r = -0.018$, $p = 0.96$). Our results reveal a \emph{detection--control gap}: styles that are linearly separable in activation space are not necessarily linearly controllable, suggesting that the model's read-out and generation pathways involve nonlinear computation that prompting can access but activation addition cannot.
